\section{Related literature and the research gap}\label{sec:lit}

\paragraph{Model averaging and forecast combination}
Since \citet{BatesGranger1969}, combining forecasts has been one of the most reliable ways to improve accuracy; \citet{Timmermann2006} and \citet{Wang2023review} review five decades.
Frequentist model averaging chooses weights by minimising estimated risk---a Mallows criterion \citep{Hansen2007}, leave-one-out or $K$-fold cross-validation \citep{HansenRacine2012,ZhangLiu2023}, or forward validation for time series \citep{ZhangZhang2023}---and establishes asymptotic optimality: the chosen weights attain the lowest attainable risk.
It now covers dependent data \citep{ZhangWanZou2013}, quantile losses \citep{LuSu2015}, general losses \citep{Yu2024unified}, time-varying weights \citep{Sun2021tvma}, covariate-dependent weights \citep{Gong2026moe}, covariate shift \citep{Zhang2026shift} and prediction intervals \citep{Qu2025intervals}.
The forecast-combination puzzle, that equal weights frequently beat estimated weights, is attributed mainly to estimation error \citep{SmithWallis2009,Claeskens2016}; remedies shrink weights towards equal weights \citep{DieboldShin2019} or restrict them \citep{Zou2025weights}.
With many series, weights are estimated series by series or learned from series features \citep[FFORMA,][]{MonteroManso2020}; \citet{AiolfiTimmermann2006} cluster \emph{forecasting models} by past performance.
All of these treat each target series separately or impose one rule on all, and none groups the \emph{target units}.

\paragraph{Decision-focused learning}
The loss that matters is the decision cost \citep{Granger1969}.
The data-driven newsvendor literature learns order quantities directly from data \citep{Levi2015,BanRudin2019}, prescriptive analytics maps covariates to decisions \citep{BertsimasKallus2020}, and ``smart predict-then-optimise'' trains predictors on decision regret \citep{ElmachtoubGrigas2022}; see \citet{Mandi2024} for a survey.
In energy forecasting, \citet{Wang2019cqra} combine probabilistic load forecasts by minimising pinball loss on the simplex, one series at a time.
Most of this work concerns single models or one unit; the exception is data pooling \citep{GuptaKallus2022}, which shrinks the solutions of many small newsvendor-type problems towards a common anchor, which pays even for unrelated problems; decision-focused shrinkage of weights towards pooled weights, one of our benchmarks and the $G=1$ case of soft GDMA, applies the same idea to forecast combination.

\paragraph{Deep learning for forecasting and for combination weights}
Pretrained foundation models such as Chronos \citep{Ansari2024chronos} forecast unseen series zero-shot with accuracy close to trained models, and deep networks increasingly \emph{produce} combination weights: FFORMA maps series features to weights with gradient boosting \citep{MonteroManso2020}, and mixture-of-experts model averaging uses a softmax neural gate \citep{Gong2026moe}.
These approaches are flexible, but fit thousands of parameters, give weights that are hard to audit and, except for \citet{Gong2026moe}, lack optimality guarantees for the combined decision.
We include a foundation model as a candidate and a decision-focused neural gate as a competitor.

\paragraph{Storage siting}
Battery siting is usually decided with engineering models of the transmission network, solved by decomposition or stochastic programming \citep{Piansky2024battery,Larroux2026bess}, which take the uncertainty storage must absorb as an input.
We make that input endogenous: residual commitment risk depends on how operators combine forecasts, so the value and best location of storage depend on the forecasting policy.

\paragraph{Latent group structures}
Panel econometrics models unobserved heterogeneity through latent groups estimated by $k$-means-type algorithms \citep{LinNg2012,BonhommeManresa2015,AndoBai2016} or classifier-Lasso \citep{Su2016}.
\citet{GuVolgushev2019} extend latent groups to panel quantile regression, the closest statistical relative of GDMA: a grouped quantile regression of demand on the candidate forecasts would be an unconstrained version of our quantile-regression benchmark.
These methods group regression parameters, typically fixed effects, for inference on the parameters; GDMA groups simplex weights over existing forecasting systems, with unit-specific critical ratios, and is judged by the out-of-sample cost of its decisions.

\paragraph{The gap}
Table~\ref{tab:gap} summarises the position of GDMA.
No existing method simultaneously (a) combines a library of existing candidate rules, including machine-learning and foundation-model forecasts, (b) chooses weights by the downstream decision cost with unit-specific cost asymmetries, and (c) shares strength across many decision units through latent segments whose number is chosen from data, with guarantees on optimality and segment recovery; nor does any link the resulting operational risk to the siting of new storage.
The ingredients exist---decision-focused combination, shrinkage across problems, latent groups---and the contribution is their combination, the theory for it, and its use for a planning decision; much of the gain over the best single forecaster is already delivered by decision-focused shrinkage, and grouping adds a smaller refinement.
The gap matters most when windows are short, because systems change and old data are stale, and units are many, as in infrastructure planning.

\begin{table}[t]
\centering
\caption{Positioning of GDMA relative to existing approaches.}\label{tab:gap}
\footnotesize
\begin{tabular}{@{}p{6.2cm}cccc@{}}
\toprule
 & Combines & Decision- & Shares strength & Heterogeneous \\
Approach & candidates & focused loss & across units & weights \\
\midrule
Per-series combination / JMA / forward validation & \checkmark & & & \checkmark \\
Pooled combination & \checkmark & & \checkmark & \\
Quantile / general-loss model averaging & \checkmark & \checkmark & & \checkmark \\
Shrinkage towards equal or pooled weights & \checkmark & & partial & \checkmark \\
Data pooling in stochastic optimisation & & \checkmark & partial & \checkmark \\
Feature-based weights (FFORMA) & \checkmark & & \checkmark & \checkmark \\
Neural mixture-of-experts gate & \checkmark & (\checkmark) & \checkmark & \checkmark \\
Data-driven newsvendor, SPO & & \checkmark & \checkmark & \\
Latent-group panel (quantile) models & & (\checkmark) & \checkmark & \checkmark \\
\textbf{GDMA (this paper)} & \checkmark & \checkmark & \checkmark & \checkmark \\
\bottomrule
\end{tabular}
\end{table}
